\begin{displayquote}
	\textit{\enquote{Nul pouvoir politique sans contrôle de l'archive, sinon de la mémoire.La démocratisation effectivement se mesure toujours à ce critère essentiel : la participation et l'accès à l'archive, à sa constitution et à son interprétation.}}
	
	\par 
	Jacques Derrida\footnote{\cite{derrida-1995}. p. 11.} \\[0.2cm]
\end{displayquote}




Par cette affirmation, l'auteur nous rappelle que l'archive ne se limite pas à un simple dépôt documentaire. Celle-ci se situe bien au-delà de cette compréhension. Elle constitue le lieu où s'élabore la mémoire individuelle, collective et institutionnelle. Sa valeur ne réside donc pas seulement dans le document lui-même, mais également dans le contexte qui l’a vu naître, dans les relations qu’il entretient avec les autres documents constitutifs d’un fonds. Ce constat renvoie au principe fondamental du respect des fonds, en ce sens qu’un document ne peut être compris s’il n’est pas lié à son contexte de production, à sa provenance et aux relations organiques qui le lient aux autres documents produits dans le cadre d’une même activité.

Cette dimension contextuelle revêt une importance particulière au sein des institutions patrimoniales, dont la mission ne se réduit pas seulement à conserver des objets, mais également à préserver les connaissances qui permettent de les comprendre, de les documenter et de les transmettre. Le musée apparaît en ce sens comme un espace de recherche, de production et de diffusion des savoirs, bien au-delà de sa fonction d’exposition des collections. L’exercice de ces missions scientifiques s’accompagne d’une importante production documentaire, qui constitue le support de ses activités de conservation, d’étude et d’exposition\footnote{\cite{maze-2017}}. Cette perception trouve sa confirmation dans la définition adoptée par le Conseil international des musées (ICOM) en 2022, qui présente le musée comme \textit{\enquote{une institution permanente, à but non lucratif et au service de la société, qui se consacre à la recherche, la collecte, la conservation, l’interprétation et l’exposition du patrimoine matériel et immatériel}}\footcite{icom-definition-2022}. Cette formulation met en lumière le fait que toute activité muséale s’accompagne d’une production documentaire qui répond à des finalités administratives, scientifiques, techniques. Elle assure la traçabilité des collections, ainsi que des activités menées par l’institution\footcite{barbant-et-al-2014}.

Parmi les institutions patrimoniales françaises, les Arts décoratifs (MAD) occupe une place singulière grâce à son histoire, à l’ampleur de ses collections et à la diversité de ses missions. Héritier de l’Union centrale des beaux-Arts appliquées à l’Industrie, fondée à la fin du XIXe siècle\footnote{\cite{brunhammer-1992}. p. 11.}, l'Union centrale des arts décoratifs (UCAD), né dans un contexte marqué par les profondes transformations industrielles de la seconde moitié du XIXe siècle\footnote{\cite{brunhammer-1992}. p. 11.}. Ses fondateurs poursuivent alors une ambition clairement affirmée : favoriser le rapprochement entre la création artistique et la production industrielle afin de promouvoir, selon leurs termes, \enquote{la réalisation du beau dans l’utile}\footnote{\cite{brunhammer-1992}. p. 11.}. Cette vocation dépasse très tôt la seule constitution de collections muséales. 

En s'installant début du XXe siècle au Pavillon de Marsan du Palais du Louvre, le musée des Arts décoratifs est créé. Il conserve aujourd’hui l’une des plus importantes collections d’arts décoratifs au monde. Elles réunissent plus d’un million et demi d’œuvres et artefacts du Moyen Âge à nos jours — mobilier, art de la table, design, mode et textile, bijoux, papiers peints, objets d’arts, verre, arts asiatiques, jouets, publicité, dessins, ou encore les photographies\footnote{\cite{mad-collections}}. Ses collections témoignent d’une recherche permanente d’harmonie entre le beau et l’utile et sont enrichies chaque année de très nombreux dons, achats et legs\footnote{\cite{mad-collections}}. À travers ses expositions temporaires, permanentes, ses publications scientifiques, ses programmes de recherche, ses actions de médiation culturelle et ses partenariats institutionnels, le MAD produit quotidiennement une documentation abondante qui constitue elle-même un patrimoine à préserver.

L'exposition occupe une place centrale dans l'activité du MAD. Elle constitue le principal dispositif par lequel l’institution construit un discours scientifique et culturel à destination du public\footnote{\cite{drouguet-gob-2016}. p. 122-123.}. Chaque projet d’exposition mobilise simultanément un ensemble de compétences scientifiques, administratives, techniques, juridiques, logistiques et financières dans lesquelles interviennent plusieurs acteurs et d'autres directions : conservateurs, régisseurs, graphistes, scénographes, restaurateurs, juristes, communicateurs, etc. Cette collaboration génère une masse importante de documents qui accompagnent l’ensemble du projet, depuis les premières réflexions jusqu’au bilan final de l’exposition. Ces documents constituent un témoignage essentiel du fonctionnement de l’institution, des choix scientifiques opérés et  des décisions administratives. Leur conservation participe pleinement à la constitution de la mémoire institutionnelle, à la compréhension des activités transversales de ses services et de la politique de gestion et valorisation des collections nationales dans la durée.

Cette production documentaire connaît depuis une vingtaine d'années une profonde mutation sous l’effet de la transition numérique. Plus récemment, l’essor des outils collaboratifs, la dématérialisation des procédures administratives, la généralisation des échanges électroniques et le développement des espaces de stockage partagés ont transformé les modalités de création, de circulation et de conservation des documents produits par les institutions culturelles. À ce titre, les dossiers d’exposition autrefois constitués de documents papiers sont désormais élaborés presque exclusivement sous forme numérique. Si le numérique joue le rôle de facilitateur dans le travail collaboratif entre les services et les partenaires extérieurs, il s’accompagne néanmoins d’une augmentation considérable du volume documentaire produit et conservé.

Une telle évolution ne constitue pas une simple substitution d’un support à un autre. Comme le souligne Pascal Robert, la dématérialisation ne correspond pas à faire disparaître la matérialité de l’information, elle transforme les modalités de sa matérialisation et de son accès\footnote{\cite{robert-2004}. p. 57.}. Avec la duplication des fichiers, la génération progressive des données triviales, la multiplication des versions des fichiers et l’absence des règles communes de classement conduisent à une accumulation documentaire dont la maîtrise devient complexe. Aujourd’hui, plus d’un organisme public ou privé se retrouve confronté à une croissance continue de ses données numériques. 

Les institutions patrimoniales ne font pas exception. La multiplication des documents bureautiques, de plans scénographiques, de photographies, de courriers, représente un volume important. L’organisation intellectuelle tend progressivement à s’effacer derrière leur seule accumulation matérielle. Lorsque les principes archivistiques ne sont pas appliqués de manière homogène, les documents demeurent présents, mais deviennent difficiles à retrouver, et de ce fait inexploitables. La difficulté ne réside donc pas dans la quantité de documents à conserver, mais plutôt dans la perte progressive des liens intellectuelles qui donnent sens à ces derniers. 

C’est dans cette perspective que s’inscrit le présent travail de recherche, réalisé dans le cadre de notre alternance au sein du département Bibliothèque – Archives – Documentation du musée des Arts décoratifs. L’évaluation des espaces de stockage utilisés pour conserver les dossiers numériques d’exposition met en évidence l’existence d’un ensemble documentaire caractérisé entre autres par une forte hétérogénéité des arborescences, la coexistence de doublons, la présence de fichiers orphelins et d’anomalies. Les chapitres suivants détaillent les différentes observations à partir d’un audit technique, qui suscite une interrogation sur la manière dont un ensemble numérique peut progressivement perdre son intelligibilité archivistique.

L’analyse de ce contexte conduit à s’interroger sur les conditions dans lesquelles les principes archivistiques peuvent être mobilisés pour répondre aux défis posés par la croissance des archives numériques. La maîtrise intellectuelle d’un ensemble documentaire ne repose pas seulement sur l’identification et la localisation matérielle des fichiers, elle suppose également de documenter les contextes dans lesquels les archives ont été créées, maintenues et utilisées. Le Conseil international des archives souligne ainsi que les métadonnées élémentaires attachées aux fichiers et aux répertoires doivent être complétées par une description des contextes afin de permettre l’identification, la compréhension, la recherche et l’utilisation des documents\footnote{\cite{ric-fad-2023}. p. 7.}. Les outils d’audit peuvent rendre visible l’arborescence, les formats, les volumes ou certaines anomalies d’un ensemble documentaire, mais ils ne suffisent pas, à eux seuls, de restituer les relations intellectuelles entre les documents, leurs producteurs, les processus dont ils résultent et les activités qu’ils retracent\footnote{\cite{ric-fad-2023}. p. 4-16.}. Dès lors, la question n’est plus seulement de conserver des fichiers, mais de préserver les informations et les liens organiques qui leur confèrent leur intelligibilité. Le modèle de référence pour les systèmes ouverts d’archivage d’information OAIS (\textit{Open Archival Information System}) rappelle à ce titre que la préservation numérique ne peut être envisagée sous un angle exclusivement technique et qu’elle doit également prendre en compte la provenance, l’authenticité et la capacité des utilisateurs futurs à interpréter l’information conservée\footnote{\cite{oais-2024}. p. 1-3.}.

 Il est apparu stratégique de prendre comme périmètre une production documentaire transversale résultant d’une activité reconnue d'utilité publique et dont la gestion est une obligation légale. Les dossiers d’exposition produits dans le cadre des activités scientifiques du MAD dont l'arriéré est conservé sur l’un de ses serveurs en font un corpus idéal. Ce travail cherche à comprendre comment un ensemble documentaire dépourvu d'organisation peut retrouver une structure conforme aux exigences archivistiques grâce à l’articulation entre expertise humaine et traitement automatisé. Mieux, l’objectif est de démontrer comment un ensemble de données en masse peut se transformer en un ensemble structuré et cohérent. Tout ceci soulève une interrogation : dans quelle mesure l’articulation entre l’audit technique, l’expertise archivistique et le traitement algorithmique permet-elle de rétablir l’intelligibilité d’un ensemble de dossiers numériques d’exposition et d’en préserver la valeur patrimoniale ? Pour répondre à cette question, trois hypothèses principales sont évoquées.
 
 La première hypothèse suppose que les méthodes traditionnelles de traitement archivistique fondées sur une analyse exclusivement manuelle atteignent leurs limites lorsqu’elles sont appliquées à des ensembles numériques caractérisés par une volumétrie importante et par la répétition d’opérations techniques. Dans ce contexte, l’automatisation ne se substitue pas à l’archiviste ; elle constitue un moyen d’exécuter de façon systématique des règles préalablement définies et validées par celui-ci. Cette hypothèse est confrontée au temps de traitement, au volume du corpus et aux résultats obtenus lors de l’expérimentation.
 
 La deuxième hypothèse part du principe que les seules propriétés techniques propres à chaque fichier, telles que l’extension, le format, la taille ou la date de modification ne suffisent pas à déterminer leur fonction (leur rôle dans l’activité qui les a générés) et leur valeur archivistiques (l’intérêt de leur conservation). L’identification automatisée des formats apporte des renseignements utiles à la caractérisation et à la préservation d’un corpus. Certains outils permettent notamment de relever les formats, leurs versions, la taille des fichiers, ainsi que la génération de métadonnées susceptibles d’aider à repérer des doublons\footnote{\cite{national-archives-droid}. p. 1.}. Toutefois, ces informations ne permettent pas, à elles seules, d'identifier l’activité à laquelle un document se rattache, ni d'en déterminer le producteur ou les liens qu’il entretient avec les autres pièces du dossier. La restitution du contexte exige donc de croiser ces propriétés techniques avec des critères archivistiques et sémantiques portant notamment sur la fonction et typologie du document, son emplacement et ses relations avec les autres composants du dossier\footnote{\cite{ric-fad-2023}. p. 7-8.}, la connaissance de l'activité productrice et du vocabulaire métier.
 
 La troisième hypothèse propose le recours à des outils de programmation, conçus à partir des besoins précis de l’institution, peut fournir une solution adaptée à certaines opérations de traitement. Cette hypothèse ne consiste pas à affirmer qu’un script développé par le langage de programmation peut remplacer un système d’archivage électronique, dont les fonctions et les finalités sont différentes. Elle vise à déterminer si un prototype local peut assister des opérations préparatoires à un versement dans un système d’archivage électronique telles que l’extraction des archives compressées, la détection de critères de tri, la proposition d’un classement et la production de journaux de traitement.
 
 Afin de vérifier ces hypothèses, notre recherche s’appuie sur une méthodologie combinant plusieurs approches complémentaires et successives. Elle repose d’abord sur une observation participante menée au sein du Département Bibliothèque-Archives-Documentation du musée des Arts décoratifs. Cette immersion professionnelle est complétée par l’étude d’un corpus de dossiers d’exposition numériques, dont le périmètre, la volumétrie et les modalités de constitution sont explicités dans les passages suivants. Un audit technique est ensuite réalisé sur la base de deux outils, l'un dédié à l'analyse et au traitement des arborescences de fichiers sous la forme de représentation visuelle ; et un outil d’aide à l’identification des formats de fichiers. Enfin, la recherche inclut la conception, le développement et l’évaluation d’un prototype de traitement automatisé initialisé par un langage de programmation. L’évaluation repose sur des indicateurs explicitement définis, tels que le nombre de fichiers traités, la proportion de classements corrects, les erreurs de rattachement, les documents non reconnus et les cas nécessitant une validation humaine.
 
 Ce mémoire s’organise donc en trois parties. La première présente la production des archives d’exposition au musée des Arts décoratifs, la place du pôle Archives dans l’écosystème documentaire de l’institution et les transformations ayant conduit à la constitution du corpus étudié. La deuxième examine les fondements archivistiques et numériques mobilisés pour penser le cycle de vie, la pérennisation et la qualification des dossiers d’exposition. Enfin, la troisième retrace la méthodologie expérimentale mise en œuvre, depuis l’audit jusqu’au développement et à l’exécution du traitement algorithmique, avant d’en évaluer les résultats, les limites et les perspectives d’application au sein de l’institution.
 
 
 
 




